\section{Introduction}

Long-context language modeling is often framed as an inference-time scaling problem: if rotary position embedding (RoPE) is compressed, interpolated, or rescaled carefully enough, a model trained at length \(L_{\text{train}}\) should extrapolate to substantially longer contexts \citep{chen2023position, peng2024yarn, ding2024longrope}. This view has produced strong engineering baselines, but it leaves one design choice essentially untouched: the training-time geometric allocation of RoPE frequencies. Standard RoPE treats that allocation as fixed, even though long-range failure is ultimately mediated through frequency collisions in the learned representation.

We take the opposite starting point. Instead of asking only how to rescale positions at inference time, we ask how frequencies should be allocated during training. We formulate RoPE frequency allocation as a variational inverse problem, derive a broadband surrogate in which the stationary density satisfies a second-order ODE, and obtain a closed-form family, EVQ-cosh. In this family, geometric RoPE appears as the degenerate \(\tau=0\) limit rather than a generic optimum.

This change in viewpoint leads to two central empirical findings. First, in the PE-dominant extreme-extrapolation regime, EVQ beats learnable adaptive positional encoding with zero additional parameters: on the DAPE-style \(128\!\rightarrow\!8\mathrm{K}\) setup, EVQ reaches \(333.7\) PPL while DAPE reaches \(455.3\). This is the cleanest setting in which positional structure itself is the dominant variable. Second, and more importantly, EVQ acts as the missing training-time foundation for inference-time scaling. Across six seeds, EVQ+YaRN reaches \(100\%\) retrieval at \(8\mathrm{K}\), while Geo+YaRN remains at \(61\%\) to \(65\%\); in a fair \(10\%\) mix comparison, EVQ+YaRN lowers PPL from \(82.9\) to \(70.9\) at \(8\mathrm{K}\) and from \(157.7\) to \(107.5\) at \(16\mathrm{K}\).

The paper's claim is therefore sharper than ``EVQ improves RoPE.'' Our claim is that standard RoPE underperforms because its geometric frequency allocation is the wrong training-time foundation for long-context extrapolation. Inference-time scaling can only recover what the checkpoint geometry makes available. Once that geometry is improved, the strongest existing inference-time methods become significantly more effective.

This distinction matters operationally. The dominant deployment pattern in long-context LLMs is to keep the training checkpoint fixed and layer stronger inference-time scaling on top. If the checkpoint has already collapsed too much low-frequency resolution into a geometric schedule, then inference-time scaling is trying to stretch positions on top of the wrong substrate. Our thesis is that training-time frequency design is therefore a missing systems primitive, not merely a cleaner mathematical interpretation of RoPE.

Our contributions are:
\begin{itemize}
\item We derive a closed-form variational solution for RoPE frequency allocation and show that geometric RoPE is the \(\tau=0\) degenerate case of a broader optimal family.
\item We show that closed-form EVQ beats learnable PE in DAPE-style extreme extrapolation, supporting the claim that PE quality itself can determine long-range behavior.
\item We demonstrate that the strongest systems result is not merely EVQ over geometric RoPE, but EVQ unlocking YaRN: training-time and inference-time positional optimization are orthogonal and strongly complementary.
\item We present a conservative empirical package in which the main claims rely on multi-seed text evidence, while larger-scale single-seed runs and cross-modal transfer results are treated as consistency checks rather than primary anchors.
\end{itemize}
